\section{References}
\label{sec:mealy-bib}

Mealy machines were introduced in \cite[Section 2.1]{mealy1955}. A related model is Moore machines~\cite[p.133]{moore1956}, in which the $n$-th output letter depends on the first $n-1$ input letters.

The Krohn-Rhodes Theorem  was originally proved in 
    \cite[p. 454]{krohnRhodes1965}. Let us remark on the differences between the original statement and the present text. In the original, semigroup terminology is used, while here we use the language of Mealy machines. The original also has an extra kind of product -- parallel composition -- instead of using only sequential composition as in the present text. Finally, the original statement says that the groups in the decomposition must have been present in the original machine, which is also true here but not stated explicitly (however, we see parts of this in the results on counter-free machine, which talk about machines that have no groups at all). The proof in this book is inspired by a proof about \textsc{ltl} from \cite[p.35]{wilke1999}.   